\subsection{Parallel Traversal}

Sometimes a program must traverse more than one iterable at the same time. One list may contain names, another may contain scores, and another may contain labels. The task is not to finish one list and then start the next. The task is to take the first value from each source, then the second value from each source, then the third value from each source, and so on.

Python provides \texttt{zip()} for this pattern.

\begin{quote}
\texttt{zip()} synchronizes several traversal sources and produces one tuple per parallel step.
\end{quote}

This is not the same thing as merging containers permanently. Like \texttt{enumerate()}, \texttt{zip()} creates a lazy iterator. It produces combined values only when the loop or \texttt{next()} asks for them.

\subsubsection{\texttt{zip()} as a Lazy Synchronization Mechanism Across Multiple Iterables}

The \texttt{zip()} function receives two or more iterables and returns a \texttt{zip} object. Each produced item is a tuple containing one value from each input iterable.

%<<<PROGRAM:programs/par08>>>
The \texttt{zip} object is an iterator. It has \texttt{\_\_iter\_\_()} and \texttt{\_\_next\_\_()}. Therefore it is consumed as it advances.

Manual inspection with \texttt{next()} shows the produced tuples:

%<<<PROGRAM:programs/par07>>>
Each produced tuple contains one value from each input iterable. The first tuple contains the first name and the first score. The second tuple contains the second name and the second score.

The tuple can be unpacked directly in the loop header:

\begin{verbatim}
names = ["Jerry", "Elaine", "Kramer"]
scores = [72, 91, 42]

for name, score in zip(names, scores):
    print(f"{name:<8} | {score:>3}")
\end{verbatim}

Possible output:

\begin{verbatim}
Jerry    |  72
Elaine   |  91
Kramer   |  42
\end{verbatim}

The loop header:

\begin{verbatim}
for name, score in zip(names, scores):
\end{verbatim}

means that each tuple produced by \texttt{zip()} is decomposed into two names. The first tuple component is assigned to \texttt{name}. The second tuple component is assigned to \texttt{score}.

This is equivalent to unpacking manually:

%<<<PROGRAM:programs/par06>>>
The unpacked version is usually clearer because the structure of each produced tuple is visible in the loop header.

\texttt{zip()} can synchronize more than two iterables:

\begin{verbatim}
names = ["Jerry", "Elaine", "Kramer"]
scores = [72, 91, 42]
labels = ["ok", "excellent", "answer"]

for name, score, label in zip(names, scores, labels):
    print(f"{name:<8} | {score:>3} | {label}")
\end{verbatim}

Possible output:

\begin{verbatim}
Jerry    |  72 | ok
Elaine   |  91 | excellent
Kramer   |  42 | answer
\end{verbatim}

Now each produced tuple has three components.

The important point is that \texttt{zip()} advances all input iterators together. One output tuple is produced from one synchronized step.

\begin{verbatim}
names:  Jerry     Elaine     Kramer
scores: 72        91         42

zip output:
        ('Jerry', 72)
        ('Elaine', 91)
        ('Kramer', 42)
\end{verbatim}

By default, \texttt{zip()} stops when the shortest input iterable is exhausted.

\begin{verbatim}
names = ["Jerry", "Elaine", "Kramer", "Newman"]
scores = [72, 91]

for name, score in zip(names, scores):
    print(f"{name:<8} | {score:>3}")

print("loop finished")
\end{verbatim}

Possible output:

\begin{verbatim}
Jerry    |  72
Elaine   |  91
loop finished
\end{verbatim}

The values \texttt{"Kramer"} and \texttt{"Newman"} are not produced because the \texttt{scores} list ends first. The default \texttt{zip()} rule is shortest-input completion.

This behavior is useful when the shortest stream naturally defines the valid synchronized region. But it can hide mistakes when the lists were expected to have equal length.

A direct length check is simple:

%<<<PROGRAM:programs/par05>>>
In Python versions that support it, \texttt{zip()} also has a \texttt{strict=True} option. With this option, unequal input lengths raise a \texttt{ValueError} instead of silently stopping at the shortest iterable.

%<<<PROGRAM:programs/par04>>>
The exact message may vary slightly by Python version, but the meaning is stable: strict zipping requires all input iterables to finish together.

The \texttt{zip()} object is lazy. It does not build all tuples immediately.

%<<<PROGRAM:programs/par03>>>
The memory address will differ between executions. The important point is that \texttt{pairs} is not a list. It is a traversal object.

When stored results are needed, the zip object can be materialized:

\begin{verbatim}
names = ["Jerry", "Elaine", "Kramer"]
scores = [72, 91, 42]

stored_pairs = list(zip(names, scores))

print(f"stored_pairs: {stored_pairs}")
print(f"type(stored_pairs): {type(stored_pairs)}")
\end{verbatim}

Possible output:

\begin{verbatim}
stored_pairs: [('Jerry', 72), ('Elaine', 91), ('Kramer', 42)]
type(stored_pairs): <class 'list'>
\end{verbatim}

The materialized result is a list of tuples. The list stores the produced pairs. Before materialization, the \texttt{zip} object produces pairs only on demand.

Like other iterators, a \texttt{zip} object is consumed after traversal:

%<<<PROGRAM:programs/par02>>>
The first \texttt{list(pairs)} consumes the zip iterator. The second one receives no remaining values.

To repeat the traversal, create a new \texttt{zip} object from the original iterables:

%<<<PROGRAM:programs/par01>>>
The original lists are reusable. Each call to \texttt{zip(names, scores)} creates a new iterator.

\texttt{zip()} can also synchronize different iterable types:

\begin{verbatim}
letters = "abc"
nums = range(1, 4)
words = ("soup", "pretzels", "muffin")

for letter, num, word in zip(letters, nums, words):
    print(f"{letter} | {num:>2} | {word}")
\end{verbatim}

Possible output:

\begin{verbatim}
a |  1 | soup
b |  2 | pretzels
c |  3 | muffin
\end{verbatim}

The string supplies characters. The \texttt{range} supplies computed integers. The tuple supplies stored references. The \texttt{zip} object only coordinates the parallel traversal.

The practical rule is:

\begin{quote}
Use \texttt{zip()} when several iterables must be consumed in synchronized steps. Remember that the result is a lazy iterator and, by default, traversal stops at the shortest input.
\end{quote}